\section{Introduction}
\label{sec:intro}

% FigTeaser.tex -- Figure 1 (the teaser): the apparatus on top, the ten repositories below.
% Source: fig_teaser_editable_2026-08-20.pdf, cropped to its ink box.
% MUST stay the first float so it numbers as Figure 1.
\begin{figure*}[t]
  \centering
  \includegraphics[width=\linewidth]{fig_teaser}
  \caption{\textbf{\ai{} at a glance.} \emph{Top:} the lifecycle every task runs through, drawn here
  with discrete graph diffusion as the running example. The repository, the model it starts from and
  the cheap proxy metric are frozen; the agent then has four hours on one B300 to read the code,
  change it, and test each idea against that proxy. What it leaves behind is a source-code patch and
  nothing else --- no weights it trained, no cached state. The patch is applied in a fresh container
  and run from initialization for up to twelve hours, and at most the three most recent checkpoints
  are kept. An evaluator fixed before the first run, with no access to the agent's workspace, scores
  them, and the cell takes the best of the three under the task's direction; for the two tasks that
  do not train, the twelve hours are only a ceiling on the verification stage
  (\S\ref{sec:protocol}). \emph{Bottom:} the ten frozen repositories, chosen so that between them
  they cover ten distinct families of training algorithm, each shown with the final metric its own
  evaluator computes and the asset that metric reads (Table~\ref{tab:tasks}).}
  \label{fig:teaser}
\end{figure*}


The premise that AI systems will soon do AI research is no longer a forecast held at a distance: it is
the stated direction of frontier laboratories and the organizing question of a fast-growing literature
on recursive self-improvement, which takes machine-learning engineering as its executable testbed
\citep{yang2026frontis,meng2026rsibench,lu2026metaagent}. The evidence base under that premise is
thinner than the attention it receives, and it is thin in a particular way. Coding agents can now open
a real repository, edit it and satisfy an executable test \citep{jimenez2024swebench}, and the
benchmarks that extend this to research measure almost exclusively whether an agent can raise a fixed,
narrow, automatically verifiable metric --- Kaggle competitions scored against public leaderboards
\citep{chan2024mlebench,qiang2025mledojo}, ML research competitions scored by the fraction of a
baseline-to-human gap they close \citep{zhang2025mlrcbench}, post-training runs scored against a
released instruction-tuned checkpoint \citep{rank2026posttrainbench}, and research-engineering
environments scored against expert attempts \citep{wijk2024rebench}. The sharpest criticism of that
literature is that the parts of research which actually decide whether a programme succeeds are not of
that form \citep{kirgis2026openended}: choosing which hypothesis to commit to,
deciding what evidence would settle it, and recognising that an approach has failed and the right move
is to start over. Those are judgements, and a leaderboard does not contain them.

This paper argues that they are nonetheless measurable, and that measuring them does not require
abandoning verifiable tasks --- it requires keeping the record of how each one was attempted. We build
an apparatus in which a coding agent is asked to improve a real research repository's \emph{own}
training recipe, and in which the resulting claim is hard to fake: the improvement is extracted as a
source-only patch, retrained from a pinned commit in a container that carries nothing forward from
exploration, and scored by an evaluator the agent cannot run. That gives an operational form of the
narrow question --- \emph{can a coding agent make a reproducible, independently verified improvement to
an existing AI system under a controlled research budget?} --- and, because every configuration also
leaves a complete audited trajectory, it gives the wider one: \emph{when an agent does this work, what
decides how it goes?}

The narrow question is worth stating carefully because it is harder to answer than it looks. A software
task has a test that defines correctness; an AI-improvement task has a metric that only \emph{proxies}
for it, and raising a proxy is not the same as doing the work the proxy stands for --- a gap that
agents are already documented to exploit, by training on the evaluation set, by substituting a
downloaded checkpoint for one they trained \citep{rank2026posttrainbench}, or by building to whatever
oracle is in the loop \citep{ma2026buildingtotest,gabor2025evilgenie}. Whether a change
counts as an improvement also depends on what it is compared against, and a research repository supplies
more than one defensible reference: the checkpoint it ships, the recipe it commits to, and that recipe
run under a different budget.

\ai{} makes that question answerable. It freezes ten research repositories, spanning generative
modelling, language-model post-training, reward modelling, compression and agentic RL, and asks an
agent to improve each repository's \emph{own} training recipe rather than to beat a synthetic target.
One apparatus serves every task and is the paper's central object (Figure~\ref{fig:pipeline}): four
hours on one GPU to read, edit, train and query a fast proxy; then a \emph{source-only} patch, a fresh
twelve-hour replay from a pinned commit, and a score from an evaluator the agent cannot run
(\S\ref{sec:benchmark}). That lifecycle is what makes a reported gain something another party can
reproduce rather than a number the agent asserts. We run it over six systems --- treating each
\emph{(model, native harness, reasoning effort)} triple as the unit under test, since a model is
inseparable from the harness that runs it --- for 290 configurations and 792 independently scored
artifacts (\S\ref{sec:setup}).

The capability result is real but modest (Figure~\ref{fig:teaser}). Agents beat the recipe the
repository itself commits to in $62.7\%$ of scored configurations, and on the three tasks whose
committed recipe is a separately tuned configuration they rarely win at all; replaying that recipe at
the agent's own budget --- an agent-free control --- leaves the rate where it was, so the margin is not
bought by budget alone (\S\ref{sec:results}). Agents reliably
\emph{recover} a competent default and only occasionally \emph{exceed} one. Claude Opus~5 leads on these runs, which we report as a standing rather than a
proven ranking: the systems below it are separated by a task or two, every cell is a single run, and
each model is inseparable from the harness that drives it.

Two things turn out to be as much of the contribution as the capability number.
The first is what the agents did with their four hours. Read as a record of behaviour rather than of
outcomes, the runs separate on something a score cannot show: \emph{where each decides the problem
lies} --- a diagnosis taken in the opening minutes from reading the repository, which on the two tasks
where the frames visibly diverge tracks the artifact more closely than the effort setting does. Where
a repository admits an unintended route, nearly every run finds it, and what separates the artifacts
is what each decided that route was for. Reasoning effort moves all of this and moves the score not at
all (\S\ref{sec:behavior}). The second is the apparatus: a best-of-three checkpoint rule alone is
worth $7.8$ points of the pooled win rate, fewer than half the resolvable wins clear twice the
evaluator's own error, and on two tasks the metric changes which artifact a benchmark would report
(\S\ref{sec:analysis}).

Five things follow: a benchmark in which every reference is measured under one frozen evaluator and
the object of improvement is a repository's own recipe (\S\ref{sec:benchmark}); a capability result
over 290 configurations, bounded by what the artifacts support (\S\ref{sec:results}); a reading of the
same runs as a record of behaviour (\S\ref{sec:behavior}); an account of how far the apparatus itself
moves the number (\S\ref{sec:analysis}); and the release of that record, one audited account per
configuration alongside every artifact and receipt (Appendix~\ref{app:traj} reads thirty in full).

% FigPipeline.tex -- AI4AI-Bench apparatus figure (Figure 2).
%
% Now a fully-vector PDF generated by scripts/make_pipeline.py
% (Figures/generated/fig_pipeline.pdf), replacing the previous 625-line TikZ
% version. That version is preserved in git history and in
% Figures/FigPipeline_tikz.tex.bak; to revert, restore it over this file.
%
% The caption string and \label{fig:pipeline} are unchanged from the TikZ
% version. MUST follow FigTeaser (Figure 1) so it numbers as Figure 2.
\providecommand{\PipCaption}{%
  One apparatus, three lanes, one example task throughout. The starting
  checkpoint and the repository's own recipe are measured rather than assumed,
  and the recipe twice: one configuration, run once at the step count it commits
  to and once at the agent's own twelve hours. Every lane enters from the same
  pinned commit and leaves through the same frozen evaluator. The agent queries
  the fast metric without limit and never sees the final metric that decides its
  score.}

\begin{figure*}[t]
  \centering
  \includegraphics[width=\linewidth]{fig_pipeline}
  \caption{\PipCaption}
  \label{fig:pipeline}
\end{figure*}

